\subsection*{MT-M34. Finite cutoffs erase a genuine Euler condition;
the exact cost of the escaping correction}

This result concerns the literal cutoff-strictness assertion in
Masullo's Theorem 11.1.2.1(2), original PDF p.~480.
It remains a counterexample after replacing every actual
Euler-difference range by its Hilbert closure as in MT-M31.
The uncut Euler quotient and its local trace survive, as
proved in MT-M32; the finite cutoff does not preserve that
quotient.

\paragraph{Statement.}
Fix $s>1$ and $h>0$, and retain the notation
\[
 w_s(x)=(1+x^2)^{s/2},\qquad
 H_s=L^2(\mathbf R,w_s(x)\,dx),\qquad
 U_hf(x)=f(x-h).
 \tag{M34.1}
\]
Let
\[
 V=\overline{(1-U_h)H_s}^{\,H_s},\qquad
 (Rf)(u)=\sum_{n\in\mathbf Z}f(u+nh),\quad0\le u<h.
 \tag{M34.2}
\]
The periodization is a bounded surjection onto $L^2(0,h)$,
and $V=\ker R$ is a proper closed subspace, by the exact
weighted quotient model (M32.14a)--(M32.14c).
If a measurable set $\Omega\subset\mathbf R$ has a
complement containing an interval of length $h$, up to
sets of measure zero, then its indicator multiplier satisfies
\[
               M_\Omega V=M_\Omega H_s.
 \tag{M34.3}
\]
This is equality of the actual images, not merely equality
of their closures.  In particular it holds for every $L>0$
and both
\[
 \Omega_{+,L}=[-L,\infty),\qquad
 \Omega_{0,L}=[-L,L].
 \tag{M34.4}
\]
The corresponding projections converge strongly to the
identity as $L\to\infty$, although $V\ne H_s$.

More generally, if $\rho\in L^\infty(\mathbf R)$ vanishes
almost everywhere on an interval of length $h$, then
\[
                   M_\rho V=M_\rho H_s.
 \tag{M34.5}
\]
This includes every bounded compactly supported cutoff.
In (M34.5) the common image need not be closed.

There is also an exact quantitative form of the obstruction.
Let $g\in H_s$ be compactly supported, put $F=Rg$, and
suppose $F\ne0$.  For either set in (M34.4), take $L$
large enough that $g$ is supported in $\Omega$.  Write
$E=\mathbf R\setminus\Omega$ and
\[
 W_E(u)=\sum_{\substack{n\in\mathbf Z\\u+nh\in E}}
                   \frac1{w_s(u+nh)}.
 \tag{M34.6}
\]
Then the minimum norm of a lift in the proper closed
subspace $V$ with prescribed cutoff $g$ is exactly
\[
 \min_{\substack{f\in V\\M_\Omega f=g}}\|f\|_{H_s}^2
       =\|g\|_{H_s}^2+
                      \int_0^h\frac{|F(u)|^2}{W_E(u)}\,du.
 \tag{M34.7}
\]
In particular, as $L\to\infty$,
\[
 \begin{split}
 \min_{\substack{f\in V\\M_{\Omega_{+,L}}f=g}}
                 \|f\|_{H_s}^2
  &=\|g\|_{H_s}^2+
        h(s-1)L^{s-1}\|F\|_{L^2(0,h)}^2(1+o(1)),\\
 \min_{\substack{f\in V\\M_{\Omega_{0,L}}f=g}}
                 \|f\|_{H_s}^2
  &=\|g\|_{H_s}^2+
       \frac{h(s-1)}2L^{s-1}\|F\|_{L^2(0,h)}^2(1+o(1)).
 \end{split}
 \tag{M34.8}
\]
Thus the norms themselves grow as $L^{(s-1)/2}$,
with leading constants $\sqrt{h(s-1)}\,\|F\|_2$
and $\sqrt{h(s-1)/2}\,\|F\|_2$, respectively.
For the source cutoff $L=\log\Lambda$ these are exact
leading constants in powers of $\log\Lambda$.

\paragraph{Repair outside the cutoff.}
For completeness, the scalar model of MT-M32 gives
\[
 W(u)=\sum_{n\in\mathbf Z}w_s(u+nh)^{-1},\qquad
 |Rf(u)|^2\le W(u)\sum_n|f(u+nh)|^2w_s(u+nh).
 \tag{M34.9}
\]
The positive continuous periodic function $W$ is bounded
above and below.  Therefore $R:H_s\to L^2(0,h)$ is
bounded.  The lift $f(u+nh)=F(u)/(W(u)w_s(u+nh))$
shows surjectivity.  The exact closure identity $V=\ker R$
is the already proved one in (M32.14c), including all
continuous Fourier coefficient constraints.  This also
shows that $V$ is proper.

Choose a half-open interval $I_0=[c,c+h)$ contained in
$\mathbf R\setminus\Omega$ up to a null set.  For almost
every $u\in[0,h)$ there is a unique integer $n(u)$ with
$u+n(u)h\in I_0$.  Define a lift $L_{I_0}$ by
\[
 (L_{I_0}F)(u+n(u)h)=F(u),\qquad
 L_{I_0}F=0\quad\hbox{off }I_0.
 \tag{M34.10}
\]
The piecewise translation $u\mapsto u+n(u)h$ preserves
Lebesgue measure.  Hence
\[
 \|L_{I_0}F\|_{H_s}^2
       \le\bigl(\sup_{x\in I_0}w_s(x)\bigr)\|F\|_2^2,
 \qquad R L_{I_0}F=F.
 \tag{M34.11}
\]
For a prescribed $g\in M_\Omega H_s$, the vector
\[
                        f=g-L_{I_0}Rg
 \tag{M34.12}
\]
belongs to $V$ and satisfies $M_\Omega f=g$.  This proves
(M34.3), and also supplies a bounded right inverse for each
fixed cutoff.  For (M34.5), start with any $f\in H_s$ and
replace it by $f-L_{I_0}Rf$ on an interval where $\rho=0$.
Its $M_\rho$ image is unchanged.  The reverse inclusions
in both image equalities follow from $V\subset H_s$.

For the two sets in (M34.4), the complements contain such
an interval for every finite $L$.  Their indicators increase
pointwise to one.  Dominated convergence applied to
$|f|^2w_s$ proves strong convergence of their multipliers
to the identity.  This proves all the qualitative assertions.

\paragraph{The minimum escaping correction.}
For a feasible $f$ in (M34.7), put $k=f-g$.  Then $k$ is
supported in $E$, $Rk=-F$, and the disjoint supports give
\[
                       \|f\|_{H_s}^2=\|g\|_{H_s}^2+
                                                    \|k\|_{H_s}^2.
 \tag{M34.13}
\]
For each fixed $L$, the measurable function $W_E$ is
positive and bounded above and below by positive constants
on $[0,h)$: the upper bound follows from $W_E\le W$;
for a lower bound, choose one full length-$h$ interval
in $E$ and use its bounded weight.  No continuity across
the endpoints modulo $h$ is needed.

Fiberwise Cauchy--Schwarz gives
\[
 |F(u)|^2\le W_E(u)
         \sum_{u+nh\in E}|k(u+nh)|^2w_s(u+nh).
 \tag{M34.14}
\]
After integration this is the lower bound in (M34.7).
Equality is attained by the measurable function
\[
 k(u+nh)=
 \begin{cases}
 -\dfrac{F(u)}{W_E(u)w_s(u+nh)},&u+nh\in E,\\[6pt]
 0,&u+nh\notin E.
 \end{cases}
 \tag{M34.15}
\]
The lower bound for $W_E$ proves that this function belongs
to $H_s$.  Its absolutely convergent periodization equals
$-F$ and its squared norm is the displayed integral in
(M34.7).  Equality in Cauchy--Schwarz on each fiber, or
strict convexity of the Hilbert norm, shows that this
minimum lift is unique up to equality almost everywhere.

Set $\varphi(t)=(1+t^2)^{-s/2}$ for $t\ge0$.
It is decreasing and
\[
 \int_L^\infty\varphi(t)\,dt
           =\frac{L^{1-s}}{s-1}(1+O(L^{-2})).
 \tag{M34.16}
\]
Indeed $\varphi(t)=t^{-s}(1+O(t^{-2}))$ uniformly for
$t\ge L\ge1$, and both powers can be integrated.
The decreasing-function comparison between a sum on a
lattice of step $h$ and its integral shows, uniformly in
the residue $u\in[0,h)$, that each one-sided tail satisfies
\[
 \sum_{u+nh>L}\varphi(u+nh)
   =\frac1h\int_L^\infty\varphi(t)\,dt+O(\varphi(L)),
 \tag{M34.17}
\]
with the same estimate for $u+nh<-L$ after reflection.
For clarity, the first point beyond $L$ lies between $L$
and $L+h$; comparing the decreasing sum with integrals
starting at that point incurs at most one summand, and
moving the integration endpoint back to $L$ incurs at
most $h\varphi(L)$.  These bounds are uniform in the
residue and also cover endpoint equalities.
Consequently
\[
 \begin{split}
 W_{(-\infty,-L)}(u)
    &=\frac{L^{1-s}}{h(s-1)}(1+o(1)),\\
 W_{\mathbf R\setminus[-L,L]}(u)
    &=\frac{2L^{1-s}}{h(s-1)}(1+o(1)),
 \end{split}
 \tag{M34.18}
\]
uniformly in $u$.  The uniform positive leading terms
allow reciprocation and integration against $|F|^2$.
Equations (M34.7) and (M34.18) prove (M34.8).

\paragraph{A precise missing bounded-lift hypothesis.}
Here is a useful general form of the failed inference.
Let $V$ be a closed subspace of a Hilbert space $H$, and
let orthogonal projections $E_n$ converge strongly to the
identity.  If a fixed $g\in H$ has lifts $v_n\in V$ with
\[
                 E_nv_n=E_ng,\qquad
                            \sup_n\|v_n\|<\infty,
 \tag{M34.19}
\]
then $g\in V$.  To prove this without a subsequence
argument, let $z\in V^\perp$.  Since
$\langle v_n-g,E_nz\rangle=0$,
\[
 |\langle g,z\rangle|
  =|\langle v_n-g,z\rangle|
  \le(\sup_n\|v_n\|+\|g\|)\,\|(1-E_n)z\|
          \longrightarrow0.
 \tag{M34.20}
\]
Thus $g\perp V^\perp$ and $g\in V$.  Equation (M34.8)
proves that this bounded-lift condition fails quantitatively
in the present cutoff problem whenever $Rg\ne0$.

\paragraph{Actual source specialization and exact refutation.}
Take $F=\mathbf Q$, the trivial continuous character of
$C_{\mathbf Q}$, $S=\{\infty\}$, and any rational prime
$p$.  Then $\operatorname{Pic}(\mathbf Z)=1$, the character
is unramified outside $S$, and the source index set
(11.1.2.1), p.~480, includes
\[
                         S_i=\varnothing,
                         \qquad S_{i'}=\{p\}.
 \tag{M34.21}
\]
These are distinct with $S_i\subset S_{i'}$; no place
of $S_i$ divides $p$, while a place of $S_{i'}$ does.
Choose any $s\ge2=2[\mathbf Q:\mathbf Q]$.
Thus every hypothesis of Theorem 11.1.2.1(2) is satisfied.

Normalize Haar measure on $C_{\mathbf Q}^1$ to have
mass one and use $x=\log|g|$ in its trivial sector.
The zero-line Hilbert space is exactly $H_s$.
Let $\Theta$ be the actual smooth theta image in this
sector.  MT-M24 and MT-M25 identify its closure with
$H_s$: the pole-removed completed rational zeta function
has no zero on $\Re z=0$, using the classical zero-free
boundary $\Re z=1$ and its functional equation.
This use of the boundary zero-free theorem requires no
GRH.  The source's one-prime operator is
$P_p=1-U_{-h}$ with $h=\log p$, as checked in MT-M31.
Since
\[
                    1-U_{-h}=-U_{-h}(1-U_h)
 \tag{M34.22}
\]
and $\ker R$ is translation invariant,
$\overline{P_pH_s}=V$; density and boundedness also give
$\overline{P_p\Theta}=V$.

Original Definition 9.3.3.3 and formulas
(9.3.3.4)--(9.3.3.5), p.~428, say that for $\Lambda>1$
the two cutoffs in Theorem 11.1.2.1 are exactly
$M_{\Omega_{+,L}}$ and $M_{\Omega_{0,L}}$ with
$L=\log\Lambda$.  The source defines its cutoff spaces
in (11.1.2.3), p.~480, as the Hilbert closures of the
cutoff smooth images.  For either multiplier $M_\Omega$,
boundedness and (M34.3) therefore imply
\[
 \overline{M_\Omega P_p\Theta}^{\,H_s}
       =\overline{M_\Omega V}^{\,H_s}
       =M_\Omega H_s
       =\overline{M_\Omega\Theta}^{\,H_s}.
 \tag{M34.23}
\]
The common subspace is closed because $M_\Omega$ is an
orthogonal projection.  Thus the source inclusion
(11.1.2.6) is equality for \emph{every} finite $\Lambda>1$,
for both stipulated types of cutoff.  In particular, for
the non-decreasing sequence $\Lambda_k=e^{k+1}$, no
subsequence has the strict inclusion asserted in
Theorem 11.1.2.1(2).  The uncut inclusion (11.1.2.5)
is nevertheless strict, since $V\ne H_s$.

The source's proof on p.~482, immediately after
(11.1.2.16), infers non-strictness of the uncut inclusion
from equality at each cutoff and strong convergence of
the cutoff multipliers to the identity.  Equations
(M34.3)--(M34.8) disprove exactly that inference and
identify the missing uniform control of its lifts.
The source later invokes Theorem 11.1.2.1 in the proof
of its geometric trace Theorem 11.3.2.1, pp.~495--496;
the main Theorem 12.5.1.1, p.~520, in turn invokes
Theorem 11.3.2.1.  This is a targeted dependency trace,
not a proof audit of the remaining kernel arguments.

\paragraph{Input and placement.}
The new cutoff and norm arguments above are elementary
and unconditional; the exact Euler closure and actual
theta-density identification are explicit antecedents
MT-M31--MT-M32 and MT-M24--MT-M25.  Original PDF
pp.~428, 480 and 482 were rendered and visually read
to check the definitions, full strictness hypotheses and
the disputed implication.  MT-M34 refutes this literal
intermediate assertion and proves its quantitative
replacement.  It does not assert that the conditional
Hecke-GRH endpoint $M_H$ is false, or that a different
geometric trace construction is impossible.

For cote 1, scans 79--80 and 111--112, this distinguishes
the trace of a genuine local Euler quotient from a
trace attributed to its finite radial cutoffs.  The
quotient supplies the local closed-point contribution;
the finite cutoff images alone carry no such quotient
condition.  A comparison of these traces therefore
requires quantitative control beyond strong convergence.
